% Macros simples prevues pour etre synchronisees avec MathJax plus tard.
\newcommand{\R}{\mathbb{R}}
\newcommand{\C}{\mathbb{C}}
\newcommand{\Q}{\mathbb{Q}}
\newcommand{\Z}{\mathbb{Z}}
\newcommand{\N}{\mathbb{N}}
\newcommand{\K}{\mathbb{K}}
\newcommand{\eps}{\varepsilon}
\newcommand{\vect}[1]{\overrightarrow{#1}}
\newcommand{\abs}[1]{\left|#1\right|}
\newcommand{\norme}[1]{\left\lVert#1\right\rVert}
\newcommand{\croch}[1]{\left[#1\right]}
\newcommand{\paren}[1]{\left(#1\right)}


\theme{fonctions-de-deux-variables}
\track{sup}
\name{Fonctions de deux variables}
\description{Dérivées partielles, gradients, extrema et étude locale en dimension 2.}

\questioncours{1}{
Question de cours à rédiger pour le thème Fonctions de deux variables.
}

\questioncours{2}{
Deuxième question de cours à rédiger pour le thème Fonctions de deux variables.
}

\exercice{1}{
\difficulty{1}
\title{Exercice 1}
\statement{
Énoncé de difficulté 1 à rédiger pour le thème Fonctions de deux variables.
}
\answer{
}
}

\exercice{2}{
\difficulty{1}
\title{Exercice 2}
\statement{
Deuxième énoncé de difficulté 1 à rédiger pour le thème Fonctions de deux variables.
}
\answer{
}
}

\exercice{3}{
\difficulty{2}
\title{Exercice 3}
\statement{
Énoncé de difficulté 2 à rédiger pour le thème Fonctions de deux variables.
}
\answer{
}
}

\exercice{4}{
\difficulty{2}
\title{Exercice 4}
\statement{
Deuxième énoncé de difficulté 2 à rédiger pour le thème Fonctions de deux variables.
}
\answer{
}
}

\exercice{5}{
\difficulty{3}
\title{Exercice 5}
\statement{
Énoncé de difficulté 3 à rédiger pour le thème Fonctions de deux variables.
}
\answer{
}
}

\exercice{6}{
\difficulty{3}
\title{Exercice 6}
\statement{
Deuxième énoncé de difficulté 3 à rédiger pour le thème Fonctions de deux variables.
}
\answer{
}
}
